\chapter*{Glossary}
\addcontentsline{toc}{section}{Glossary}

Plain-language definitions of every term of art the document leans on,
for reading without the original papers at hand. Terms are grouped by
the section that introduces them.

\section*{The problem (Chapter~\ref{sec:problem})}

\begin{description}[leftmargin=1.5em, itemsep=1pt, font=\normalfont\bfseries]
\item[trial] One training run with one fixed configuration and one
recorded random seed.
\item[incumbent] The best configuration found so far in a study.
\item[fidelity] The resource level of a partial evaluation: epochs
trained, environment steps consumed, chain length run. Low fidelity is
cheap and only loosely predictive of full results.
\item[searcher] The half of a tuner that decides which configuration to
try next.
\item[scheduler] The half of a tuner that decides how long trials run,
which are stopped early, and which paused ones resume.
\item[run store] The database recording every trial's configuration,
seed, curves, cost, and diagnostics; the raw material for warm starts
and transfer.
\item[seed protocol] Evaluating candidates on several training seeds and
reporting results on separate test seeds the tuner never saw, so that a
study optimizes configurations rather than luck.
\end{description}

\section*{Model-free search (Chapter~\ref{sec:modelfree})}

\begin{description}[leftmargin=1.5em, itemsep=1pt, font=\normalfont\bfseries]
\item[log-uniform sampling] Drawing a scale-type knob (like a learning
rate) uniformly in its exponent, so each factor-of-ten range gets equal
attention.
\item[Sobol points] A deterministic sequence that spreads samples more
evenly than independent random draws, avoiding accidental clusters and
gaps.
\item[effective dimensionality] The number of knobs that actually move
the objective; usually far smaller than the number of knobs, but which
ones matter varies by problem.
\item[CMA-ES] An evolutionary optimizer that keeps a Gaussian sampling
distribution and reshapes it toward directions where recent samples
scored well; needs hundreds of cheap evaluations to shine.
\item[differential evolution] An evolutionary rule that proposes new
points by adding scaled differences of existing population members;
the engine inside DEHB.
\end{description}

\section*{Bayesian optimization (Chapter~\ref{sec:bayesian})}

\begin{description}[leftmargin=1.5em, itemsep=1pt, font=\normalfont\bfseries]
\item[surrogate] A cheap statistical model of the expensive objective,
fit to the trials seen so far and consulted thousands of times between
real trials.
\item[Gaussian process (GP)] The default surrogate: a Bayesian
regression over functions whose predictions come with calibrated
uncertainty; the same mathematics as kriging in spatial statistics.
\item[kernel] The GP's assumption about how similar the objective is at
two configurations, as a function of the distance between them.
\item[lengthscale] The distance along one axis over which the objective
decorrelates under the kernel; a long fitted lengthscale means that
axis barely matters.
\item[marginal likelihood] The score maximized to fit kernel
parameters; it trades off fitting the data against model complexity
automatically.
\item[input/output warping] Transforming the search axes and the
observed losses so they look well behaved (stationary, roughly
Gaussian) to the surrogate; where HEBO's challenge-winning gains came
from.
\item[acquisition function] The rule that converts the surrogate's
beliefs into the next trial, trading off promising means against
informative uncertainty.
\item[expected improvement (EI)] The classic acquisition: the expected
amount by which a candidate beats the incumbent, an option-value
computation with the incumbent as strike.
\item[LogEI] The numerically repaired EI family that optimizes the
logarithm through stable formulas; identical optima, usable gradients.
\item[UCB] Optimism under uncertainty: score each point by the best
value its posterior plausibly allows, then pick the best score.
\item[fantasies] Posterior samples standing in for the unknown outcomes
of still-running trials, so new proposals can be made without waiting.
\item[TPE] A sampler that models where good and bad trials sit in
configuration space and proposes where the good-to-bad density ratio
peaks; native to conditional spaces, blind to interactions.
\item[trust region] A box around the incumbent inside which the
surrogate is trusted; grows after successes, shrinks after failures,
restarts on collapse.
\item[ask-and-tell] The interface shape where the tuner is asked for a
configuration and later told the result, with no control inversion.
\end{description}

\section*{Multi-fidelity (Chapter~\ref{sec:multifidelity})}

\begin{description}[leftmargin=1.5em, itemsep=1pt, font=\normalfont\bfseries]
\item[rung] One fidelity level in a race: all trials at a rung are
compared at the same budget.
\item[successive halving] Race many configurations at a small budget,
keep the best fraction, multiply their budget, repeat.
\item[Hyperband] A portfolio of successive-halving races from
aggressive to conservative, hedging the choice of starting budget.
\item[ASHA] Asynchronous successive halving: promote the moment a trial
is in the top fraction of results seen so far, never waiting for a rung
to fill; the standard distributed scheduler.
\item[median stopping rule] Kill any trial running below the median of
completed trials at the same step; the cheapest insurance.
\item[BOHB / DEHB] Racing schedulers whose new entrants come from a
model (TPE) or from differential evolution instead of random sampling;
DEHB is the best-evidenced small-budget RL tuner.
\item[freeze-thaw / learning-curve scheduling] Keeping a basket of
paused trials and allocating the next budget slice by extrapolating
each trial's curve with uncertainty.
\item[ifBO] The learning-curve scheduler whose extrapolator is a
transformer pre-trained once on synthetic curves, so no model is refit
inside the loop.
\end{description}

\section*{Population methods (Chapter~\ref{sec:population})}

\begin{description}[leftmargin=1.5em, itemsep=1pt, font=\normalfont\bfseries]
\item[PBT] Training a population whose members periodically copy the
weights and hyperparameters of better members (exploit) and perturb
them (explore); the output is a schedule, not a point.
\item[truncation selection] The exploit rule: bottom performers restart
from copies of top performers.
\item[lineage / schedule] The chain of copies and perturbations behind
the final best agent; the hyperparameter trajectory it traces is the
method's real product.
\item[PB2] PBT with the explore step replaced by a time-varying GP
bandit maximizing a UCB rule over performance increments.
\item[BG-PBT] The flagship population method: trust-region mixed-space
BO for the explore step, plus generational architecture search with
distillation to warm-start new networks.
\item[distillation] Training a new network to imitate a trained one
before switching to the real objective; what makes changing
architecture mid-run affordable.
\end{description}

\section*{Multiple objectives (Chapter~\ref{sec:moo})}

\begin{description}[leftmargin=1.5em, itemsep=1pt, font=\normalfont\bfseries]
\item[dominance / Pareto front] One configuration dominates another if
it is at least as good everywhere and better somewhere; the front is
the set of undominated trade-offs, exactly as in welfare economics.
\item[hypervolume] The area (volume) of objective space dominated by
the current front, up to a fixed reference point; the standard scalar
progress measure for front-seeking methods.
\item[scalarization] Collapsing several objectives into one number;
the Chebyshev form (weighted worst gap to the ideal) can reach every
Pareto point, the weighted sum cannot.
\item[EHVI / qNEHVI / qLogNEHVI] Expected hypervolume improvement and
its batched, noise-aware, numerically repaired successors; the premium
acquisition for two to four objectives.
\item[non-dominated sorting] Layering a population into ranks
(undominated; undominated once rank one is removed; and so on), used
for selection and for rung promotion.
\item[crowding distance] A tie-breaker preferring points in sparsely
covered stretches of the front, so selection spreads along it.
\item[veto gate] A correctness check (divergences, $\widehat{R}$) that
discards a trial regardless of its score; distinct from an objective.
\end{description}

\section*{Priors, transfer, and scale (Chapter~\ref{sec:transfer})}

\begin{description}[leftmargin=1.5em, itemsep=1pt, font=\normalfont\bfseries]
\item[$\pi$BO] Multiplying the acquisition by an expert's prior over
where the optimum lies, raised to a power that decays as evidence
arrives; steers early, forgets bad advice.
\item[PriorBand] The same idea inside multi-fidelity racing: new
configurations drawn from a portfolio of uniform, prior, and
around-the-incumbent samplers, weighted by rung.
\item[warm start] Seeding a new study with configurations that ranked
well on similar past studies, straight from the run store.
\item[prior-fitted network (PFN)] A transformer pre-trained on
synthetic datasets from a chosen prior so that, shown a study's
observations as context, it outputs posterior predictions in one
forward pass; approximate Bayesian inference paid for once, offline.
\item[EI-per-cost / cost cooling] Dividing the acquisition by
predicted trial cost, with an exponent annealed toward zero so
bargains matter early and stop mattering at the end.
\item[$\mu$P / $\mu$Transfer] A network parametrization under which
good hyperparameters become approximately width-invariant, so one
tunes a small proxy and transfers to the large model.
\item[scaling laws (for tuning)] Fitted formulas giving optimal
learning rate, batch size, or weight decay as functions of model and
data scale; at the largest scales they replace search outright.
\end{description}

\section*{Architecture search (Chapter~\ref{sec:workloads})}

\begin{description}[leftmargin=1.5em, itemsep=1pt, font=\normalfont\bfseries]
\item[one-shot NAS / supernet] Training one network containing every
candidate architecture as a weight-sharing sub-network, then scoring
candidates without retraining; economical in principle, unreliable in
the documented evaluations.
\item[zero-cost proxy] A score computed at initialization, without
training, meant to predict a trained architecture's quality; best
evidence says useful only alongside a learned predictor, with plain
parameter count and FLOPs surprisingly competitive.
\item[grammar-based search space] Architectures generated by
production rules like sentences from a grammar, compactly defining
huge hierarchical spaces; solves expressiveness, not compute.
\end{description}

\section*{Sampler diagnostics (Chapter~\ref{sec:workloads})}

\begin{description}[leftmargin=1.5em, itemsep=1pt, font=\normalfont\bfseries]
\item[ESS per gradient] Effective sample size earned per gradient
evaluation; the efficiency currency for tuning MCMC, since gradients
dominate cost.
\item[$\widehat{R}$ (potential scale reduction)] A convergence
diagnostic comparing chains; values away from one veto a
configuration no matter how efficient it looks.
\item[divergence (HMC)] A numerical failure of the integrator flagged
by the sampler; counted and used as a veto.
\end{description}
